\section{Detailed Related Work}
\label{appendix-sec:detailed_related_work}
\textbf{RL for LLM reasoning.}
Reinforcement learning has become a key technique for post-training LLMs~\cite{christiano2017deep,bai2022training}. Early RLHF work~\cite{ouyang2022training,rafailov2023direct} used human preference signals to improve instruction following and alignment. More recent RLVR systems replace learned reward models with verifiable task rewards, producing large gains in mathematics and programming~\cite{guo2025deepseek,hu2025reinforce++,liu2503understanding,jaech2024openai,Zhang2025SRPOAC}. Building on RLVR, algorithmic work has improved PPO-style and GRPO-style training. VC-PPO~\cite{yuan2025s} and VAPO~\cite{yue2025vapo} strengthen value-function learning, while RLOO~\cite{ahmadian2024back} and GRPO~\cite{shao2024deepseekmath} avoid explicit value learning with multi-sample baselines. Other studies introduce replay mechanisms~\cite{liang2025squeeze,li2025repo,zhang2025rlep,lu2025safe} or adjust optimization objectives to improve stability and reduce bias~\cite{liu2503understanding,chu2025gpg,huang2025mapo,zheng2025group,www2025TFDCon,yyt2024eccv}. These works are complementary to HIVE: they improve how selected rollouts update the policy, whereas HIVE improves which prompts receive rollouts before the update.

\textbf{Data-efficient LLM training.}
Prompt and trajectory quality strongly affect LLM training efficiency. Offline curation follows the ``less is more'' principle from supervised fine-tuning~\cite{zhou2023lima,ye2025limo,zheng2024picture,wang2025comprehensivesurveyllmagentstack,tkde2024DConRec} and reinforced fine-tuning~\cite{Wang2025ReinforcementLF,Fatemi2025ConciseRV,Li2025LIMRLI,shi2025efficient,tang2025towards,cikm2022DcRec}. These approaches can reduce dataset size before training, but they cannot respond to the model's changing competence during RLVR.

Online selection addresses this limitation by conditioning selection on the current policy. Rollout-based methods discard uninformative prompts via real-time sampling~\cite{yu2025dapo,meng2025mm,foster2025learning,xu2025not,lin2025cppo,sun2025efficient} or prioritize medium-difficulty prompts~\cite{bae2025online}. They are online and precise, but they require repeated response generation before filtering. Estimator-based methods instead maintain historical logs~\cite{zheng2025act} or Bayesian utility estimates~\cite{chen2025self,zeng2025cures,shen2025bots}. These methods are efficient, but stale metadata can misrepresent the current policy's learning edge. Auxiliary-model approaches use external scorers or jointly trained predictors to estimate difficulty, as in DOTS~\cite{sun2025dots} and PCL~\cite{Gao2025PCL}; these improve adaptivity but add extra computation or memory. Stage-wise curriculum refreshes~\cite{zhang2025learning,Zhang2025SRPOAC} periodically re-estimate difficulty, but they still do not provide cheap per-step online verification.

HIVE is motivated by the observation that none of these categories simultaneously satisfy the online, and efficient requirements. Its history-informed stage inherits the low overhead of historical filtering, while its prompt-side entropy verifier corrects stale metadata without generating responses.
